\section{Conclusions}\label{sec:conclusion}

Phase I of the SS-PEG experimental program has established a single, extraordinary result:

\begin{quote}
\textbf{The entire algebraic structure of gauge symmetry---commutation relations, Jacobi identities, Casimir operators, dual Coxeter numbers, representation invariants---emerges as the unique output of a purely geometric variational principle acting on the Killing metric. No algebraic axioms are required.}
\end{quote}

\subsection{Principal Results}\label{sec:principal_results}

The experimental program reports results from 51 independent test programs, 17 Lie algebras, 4 classical families (SU, SO, Sp, SL), 3 exceptional algebras (G$_2$, F$_4$, E$_6$), abelian emergence, non-compact real forms, and $\sim 1650$ optimization runs. The key findings are:

\begin{enumerate}
\item \textbf{Killing-only ablation (Test 3.1):} The Killing metric alone uniquely determines the algebra. The non-commutativity tension $T_{\mathrm{nc}}$ is redundant. Removing $T_K$ and keeping only $T_{\mathrm{nc}}$ produces no algebraic structure. The Killing metric---defined through commutators as $K_{ab} = \mathrm{Tr}(\mathrm{ad}_a \circ \mathrm{ad}_b)$---contains \textit{more} information than the commutators themselves.

\item \textbf{Universality across families:} Every classical simple compact Lie algebra family emerges with 100\% success: SU($N$) for $N=2\ldots 5$, SO($N$) for $N=3\ldots 6$, Sp($2N$) for $N=1\ldots 3$. The exceptional algebras G$_2$, F$_4$, E$_6$ also emerge with 100\% success under constructive conditions. The only input that changes between families is the parameterization class.

\item \textbf{Isomorphism verification:} SU(2) $\cong$ Sp(2), SU(4) $\cong$ SO(6), and SO(5) $\cong$ Sp(4) yield bit-identical invariants from different parameterizations, proving that the variational principle accesses the abstract algebra, not a representation artifact.

\item \textbf{Density phase transition:} A mathematical step function at $\rho = n_{\mathrm{gen}}/(d^2-1) = 1$---0\% closure for $n_{\mathrm{gen}} = 2\ldots 14$, 100\% at $n_{\mathrm{gen}} = 15$ for SU(4). There are no intermediate states. No ``partial SU(4)''.

\item \textbf{Spacetime symmetry:} so(3,1) emerges with 100\% success rate (5/5 seeds) from $\eta$-preserving parameterization with indefinite target. The Killing eigenvalue signature is $(3+,3-)$---exactly 3 rotation generators and 3 boost generators. The conformal algebra $\su(2,2)$ also emerges with 100\% success.

\item \textbf{Abelian emergence:} U(1) appears in the Killing kernel when trace is permitted. With $d^2$ generators, the optimizer discovers $\mathfrak{u}(d) = \mathfrak{u}(1) \oplus \mathfrak{su}(d)$, with the abelian factor appearing as a zero eigenvalue of the Killing matrix.

\item \textbf{Universal $h^\vee$ formula:} $h^\vee = I_R \cdot |K|/\kappa^2$ relating Killing eigenvalues to the dual Coxeter number, verified across 15 algebras with zero exceptions.

\item \textbf{F$_4$ bifurcation:} The basin of F$_4$ exists in the extended phase space $(\theta, m, v)$ of parameters plus Adam momentum plus Adam second moment. With fresh Adam state: P(F$_4$) = 0\%. With preserved Adam state from a successful trajectory: P(F$_4$) = 100\%. The symmetry that crystallizes depends on the dynamical history of the system.

\item \textbf{Landscape measure:} A Monte Carlo analysis with 128 random seeds in $\mathfrak{so}(27)$ reveals that non-simple products dominate (53.1\% of seeds converge to closed but non-simple structures). Simple exceptional algebras have vanishingly small basins of attraction. The Standard Model's non-simple product structure is the generic outcome of the variational principle.

\item \textbf{Why the SM:} The Standard Model gauge group SU(3) $\times$ SU(2) $\times$ U(1) is the unique spectrally optimal outcome of the Killing metric variational principle, selected by the confluence of three independent criteria: topological dominance (non-simple products), spectral optimality (Ramanujan property), and accessibility (largest basin of attraction).

\item \textbf{Gravity from entanglement:} The Einstein equations emerge from entanglement thermodynamics $\delta S = \delta\langle K \rangle$ (Jacobson 1995). The same relational density $\rho$ that controls the gauge symmetry phase transition also controls spacetime curvature. Gauge forces and gravity emerge simultaneously at $\rho_c \approx 1$.

\item \textbf{Dark matter as connection:} The framework predicts that ``dark matter'' is not a particle but the footprint of the Levi-Civita connection in regions of intermediate relational density. Distinctive predictions (filamentary topology, cored profiles, no self-annihilation) align with recent ultrafaint dwarf galaxy observations~\cite{SanchezAlmeitaTrujilloPlastino2024} and are testable via weak lensing surveys.
\end{enumerate}

\subsection{Contributions of Paper I}\label{sec:contributions}

Paper I makes six contributions to the foundations of theoretical physics:

\begin{enumerate}
\item \textbf{Demonstration of algebra emergence without axioms.} The commutation relations, Jacobi identities, Casimir operators, and all representation-theoretic invariants emerge from a purely geometric variational principle. No algebraic axioms are imposed. This inverts a century of theoretical tradition: symmetry is not an axiom of gauge theory but an emergent consequence of relational geometry.

\item \textbf{The Killing-only ablation.} $T_K$ contains all the information needed to determine the algebra. $T_{\mathrm{nc}}$ is redundant. The Killing metric---defined through commutators---contains \textit{more} information than the commutators themselves.

\item \textbf{The density phase transition.} Algebraic closure is a step function: 0\% below $\rho_c = 1$, 100\% at $\rho_c = 1$. There is no ``partial symmetry.'' This is the gauge-theoretic analogue of percolation. Crucially, $\rho_c$ is also a \textbf{gravity threshold}: below $\rho_c$, there is no algebra and no curvature; above $\rho_c$, gauge forces and spacetime geometry emerge simultaneously.

\item \textbf{Spacetime emergence.} so(3,1) emerges from the same principle with indefinite Killing target, with signature $(3+,3-)$ exactly. Spacetime is not the arena---it is the pattern of non-compact directions in the emergent Killing form.

\item \textbf{Gravity from entanglement.} The P$\to$R$\to$C$\to$I mapping unifies gauge theory and general relativity at the structural level. The Riemann tensor arises as the cycle discrepancy of the Levi-Civita connection in the same way that the gauge field strength arises as the cycle discrepancy of the gauge connection. Following Jacobson~\cite{Jacobson1995,Jacobson1998}, the Einstein equations emerge from entanglement thermodynamics: gravity is the thermodynamic equilibrium condition of the entanglement structure in the evolution graph.

\item \textbf{Dark matter as connection without particles.} The framework predicts that ``dark matter'' is not a particle but the footprint of the Levi-Civita connection in regions of intermediate relational density. This yields distinctive predictions---filamentary/membrane topology, cored profiles, no self-annihilation---that are testable via weak lensing surveys and consistent with recent ultrafaint dwarf galaxy observations~\cite{SanchezAlmeitaTrujilloPlastino2024}.
\end{enumerate}

\subsection{Next Steps}\label{sec:next_steps}

The path forward is organized into three papers:

\begin{itemize}
\item \textbf{Paper II} (``Graph Topology \& Spectral Analysis'') develops the graph-topological interpretation, classifying edges as cyclic, free, or decoupled according to the Killing spectrum, and establishing the Ramanujan property as a criterion of spectral optimality.

\item \textbf{Paper III} (``Standard Model from Graph Topology'') derives the Standard Model mass spectrum, coupling constants, and flavor mixing from the same graph topology with sub-percent precision.

\item \textbf{Continuum limit.} All current experiments work with finite-dimensional matrix representations. The connection to continuum field theory requires taking $d \to \infty$ in a controlled way, ideally showing that the gauge field equations of motion emerge in this limit.

\item \textbf{E$_6$ blind discovery.} E$_6$ resists blind discovery from random initial conditions (0/144 trials). Higher $\gamma$ (500--1000), lr annealing Stage 3, or adjoint representation (78-dim) may enable discovery.

\item \textbf{E$_7$ and E$_8$.} These require HPC-class resources. E$_8$ is particularly challenging (248 generators, dimension 248).

\item \textbf{Dark matter phenomenology.} The filamentary/membrane topology prediction requires dedicated weak lensing analysis. Upcoming surveys (Euclid, Rubin, DESI) have the sensitivity to distinguish connection-based dark matter from CDM. A dedicated observational program is needed to test the void-correlation and substructure predictions.

\item \textbf{Quantitative $\rho \to$ metric mapping.} The entanglement-to-gravity chain (Section~\ref{sec:gravity_entanglement}) provides the qualitative mechanism. A quantitative mapping from $\rho(x)$ to $g_{\mu\nu}(x)$ in the continuum limit would make the gravity emergence fully predictive.
\end{itemize}
